% Kalman (1960) 图 1–4 共享样式（原生适配 html-tex-doc 深浅色自适应规范色）。
% 用法：\documentclass[tikz,border=4pt]{standalone} 后 \input{common} 再 % Kalman (1960) 图 1–4 共享样式（原生适配 html-tex-doc 深浅色自适应规范色）。
% 用法：\documentclass[tikz,border=4pt]{standalone} 后 \input{common} 再 % Kalman (1960) 图 1–4 共享样式（原生适配 html-tex-doc 深浅色自适应规范色）。
% 用法：\documentclass[tikz,border=4pt]{standalone} 后 \input{common} 再 % Kalman (1960) 图 1–4 共享样式（原生适配 html-tex-doc 深浅色自适应规范色）。
% 用法：\documentclass[tikz,border=4pt]{standalone} 后 \input{common} 再 \input{kalman-style}。
\usepackage{bm,pifont}
\usetikzlibrary{positioning,fit,calc,arrows.meta,backgrounds,shapes.geometric}

% 最优解空心亮星上标（与正文 ✫ 一致）
\newcommand{\optstar}{\text{\ding{75}}}

% 深浅色自适应规范色（配合 htd/svgtheme.py）：
% black -> currentColor；black!6 -> currentColor (opacity≈0.059)；
% white -> var(--bg)；accent (#1F6FB2) -> var(--accent)。
\colorlet{edge}{black}
\colorlet{blk}{black!6}

\tikzset{
  kalman/.style={font=\small, >={Stealth[length=2.4mm]}, line width=0.7pt, draw=edge},
  block/.style={draw=edge, fill=blk, rounded corners=1.5pt,
                minimum height=8mm, minimum width=13mm, inner sep=3pt, text=edge},
  sum/.style={draw=edge, fill=white, circle, minimum size=6.5mm, inner sep=0pt, text=edge},
  tri/.style={draw=edge, fill=blk, isosceles triangle, isosceles triangle apex angle=60,
              shape border rotate=0, inner sep=1pt, minimum height=13mm,
              font=\footnotesize, align=center, text=edge},
  sig/.style={->, edge},
  lbl/.style={font=\footnotesize, text=edge},
  dot/.style={circle, fill=edge, inner sep=1.1pt},
  frame/.style={draw=accent, dashed, rounded corners=3pt, line width=0.6pt},
  framelbl/.style={lbl, text=accent, font=\scriptsize, anchor=south},
}

% 求和节点两个输入侧 (左/下) 的 + 号，置于圆外
\newcommand{\sumsigns}[1]{%
  \node[font=\scriptsize, text=edge] at ($(#1.west)+(-0.12,0.16)$) {$+$};
  \node[font=\scriptsize, text=edge] at ($(#1.south)+(-0.16,-0.12)$) {$+$};}

\newcommand{\TXTdelay}{单位\\延迟}
\newcommand{\TXTmodel}{随机过程模型}
\newcommand{\TXTplant}{被控物理系统}
。
\usepackage{bm,pifont}
\usetikzlibrary{positioning,fit,calc,arrows.meta,backgrounds,shapes.geometric}

% 最优解空心亮星上标（与正文 ✫ 一致）
\newcommand{\optstar}{\text{\ding{75}}}

% 深浅色自适应规范色（配合 htd/svgtheme.py）：
% black -> currentColor；black!6 -> currentColor (opacity≈0.059)；
% white -> var(--bg)；accent (#1F6FB2) -> var(--accent)。
\colorlet{edge}{black}
\colorlet{blk}{black!6}

\tikzset{
  kalman/.style={font=\small, >={Stealth[length=2.4mm]}, line width=0.7pt, draw=edge},
  block/.style={draw=edge, fill=blk, rounded corners=1.5pt,
                minimum height=8mm, minimum width=13mm, inner sep=3pt, text=edge},
  sum/.style={draw=edge, fill=white, circle, minimum size=6.5mm, inner sep=0pt, text=edge},
  tri/.style={draw=edge, fill=blk, isosceles triangle, isosceles triangle apex angle=60,
              shape border rotate=0, inner sep=1pt, minimum height=13mm,
              font=\footnotesize, align=center, text=edge},
  sig/.style={->, edge},
  lbl/.style={font=\footnotesize, text=edge},
  dot/.style={circle, fill=edge, inner sep=1.1pt},
  frame/.style={draw=accent, dashed, rounded corners=3pt, line width=0.6pt},
  framelbl/.style={lbl, text=accent, font=\scriptsize, anchor=south},
}

% 求和节点两个输入侧 (左/下) 的 + 号，置于圆外
\newcommand{\sumsigns}[1]{%
  \node[font=\scriptsize, text=edge] at ($(#1.west)+(-0.12,0.16)$) {$+$};
  \node[font=\scriptsize, text=edge] at ($(#1.south)+(-0.16,-0.12)$) {$+$};}

\newcommand{\TXTdelay}{单位\\延迟}
\newcommand{\TXTmodel}{随机过程模型}
\newcommand{\TXTplant}{被控物理系统}
。
\usepackage{bm,pifont}
\usetikzlibrary{positioning,fit,calc,arrows.meta,backgrounds,shapes.geometric}

% 最优解空心亮星上标（与正文 ✫ 一致）
\newcommand{\optstar}{\text{\ding{75}}}

% 深浅色自适应规范色（配合 htd/svgtheme.py）：
% black -> currentColor；black!6 -> currentColor (opacity≈0.059)；
% white -> var(--bg)；accent (#1F6FB2) -> var(--accent)。
\colorlet{edge}{black}
\colorlet{blk}{black!6}

\tikzset{
  kalman/.style={font=\small, >={Stealth[length=2.4mm]}, line width=0.7pt, draw=edge},
  block/.style={draw=edge, fill=blk, rounded corners=1.5pt,
                minimum height=8mm, minimum width=13mm, inner sep=3pt, text=edge},
  sum/.style={draw=edge, fill=white, circle, minimum size=6.5mm, inner sep=0pt, text=edge},
  tri/.style={draw=edge, fill=blk, isosceles triangle, isosceles triangle apex angle=60,
              shape border rotate=0, inner sep=1pt, minimum height=13mm,
              font=\footnotesize, align=center, text=edge},
  sig/.style={->, edge},
  lbl/.style={font=\footnotesize, text=edge},
  dot/.style={circle, fill=edge, inner sep=1.1pt},
  frame/.style={draw=accent, dashed, rounded corners=3pt, line width=0.6pt},
  framelbl/.style={lbl, text=accent, font=\scriptsize, anchor=south},
}

% 求和节点两个输入侧 (左/下) 的 + 号，置于圆外
\newcommand{\sumsigns}[1]{%
  \node[font=\scriptsize, text=edge] at ($(#1.west)+(-0.12,0.16)$) {$+$};
  \node[font=\scriptsize, text=edge] at ($(#1.south)+(-0.16,-0.12)$) {$+$};}

\newcommand{\TXTdelay}{单位\\延迟}
\newcommand{\TXTmodel}{随机过程模型}
\newcommand{\TXTplant}{被控物理系统}
。
\usepackage{bm,pifont}
\usetikzlibrary{positioning,fit,calc,arrows.meta,backgrounds,shapes.geometric}

% 最优解空心亮星上标（与正文 ✫ 一致）
\newcommand{\optstar}{\text{\ding{75}}}

% 深浅色自适应规范色（配合 htd/svgtheme.py）：
% black -> currentColor；black!6 -> currentColor (opacity≈0.059)；
% white -> var(--bg)；accent (#1F6FB2) -> var(--accent)。
\colorlet{edge}{black}
\colorlet{blk}{black!6}

\tikzset{
  kalman/.style={font=\small, >={Stealth[length=2.4mm]}, line width=0.7pt, draw=edge},
  block/.style={draw=edge, fill=blk, rounded corners=1.5pt,
                minimum height=8mm, minimum width=13mm, inner sep=3pt, text=edge},
  sum/.style={draw=edge, fill=white, circle, minimum size=6.5mm, inner sep=0pt, text=edge},
  tri/.style={draw=edge, fill=blk, isosceles triangle, isosceles triangle apex angle=60,
              shape border rotate=0, inner sep=1pt, minimum height=13mm,
              font=\footnotesize, align=center, text=edge},
  sig/.style={->, edge},
  lbl/.style={font=\footnotesize, text=edge},
  dot/.style={circle, fill=edge, inner sep=1.1pt},
  frame/.style={draw=accent, dashed, rounded corners=3pt, line width=0.6pt},
  framelbl/.style={lbl, text=accent, font=\scriptsize, anchor=south},
}

% 求和节点两个输入侧 (左/下) 的 + 号，置于圆外
\newcommand{\sumsigns}[1]{%
  \node[font=\scriptsize, text=edge] at ($(#1.west)+(-0.12,0.16)$) {$+$};
  \node[font=\scriptsize, text=edge] at ($(#1.south)+(-0.16,-0.12)$) {$+$};}

\newcommand{\TXTdelay}{单位\\延迟}
\newcommand{\TXTmodel}{随机过程模型}
\newcommand{\TXTplant}{被控物理系统}
